\chapter{Reel Functions of Several Reel Variables : limits and continuity}

\section{Definitions and Examples}
\begin{definition}[Function of Several Variables]
    Let $E \subset \Rn$ be a non-empty set. A function $f: E \to \R$ is an application that assigns to each point $\bar{x} = (x_1, \ldots, x_n) \in E$ a reel number. $E$ is called the domain of $f$.
\end{definition}

\begin{eg}
    Consider the function $f(x,y) = \sqrt{1 - (x^2 + y^2)}$. The domain of $f$ is the set of points $(x,y)$ such that:
    \[
        1 - (x^2 + y^2) \geq 0 \Implies x^2 + y^2 \leq 1
    \]
    Therefore, the domain of $f$ is the closed unit disk in $\R^2$. We can also express $f$ as:
    \[
        z = \sqrt{1 - (x^2 + y^2)} \Iff \begin{cases}
            z \geq 0 \\
            z^2 = 1 - x^2 - y^2
        \end{cases} \Iff \begin{cases*}
            z \geq 0 \\
            x^2 + y^2 + z^2 = 1
        \end{cases*}
    \]
    The graph of $f$ is the upper hemisphere of the unit sphere in $\R^3$. Graphically:
  \begin{center}
    \begin{tikzpicture}
        \begin{axis}[
            axis lines=box,
            axis equal image,
            axis on top,
            z buffer=sort,
            view={35}{20},
            xlabel={$x$},
            ylabel={$y$},
            zlabel={$z$},
            xmin=-1, xmax=1,
            ymin=-1, ymax=1,
            zmin=0,  zmax=1,
            xtick={-1,-0.5,0,0.5,1},
            ytick={-1,-0.5,0,0.5,1},
            ztick={0,0.5,1},
            ticks=both,
            enlargelimits=false,
        ]

            \addplot3[
                surf,
                domain=0:360,
                y domain=0:1,
                samples=40,
                samples y=20,
                shader=interp,
                opacity=0.8,
                colormap={onlyprimary}{
                    color=(primary)
                    color=(secondary)
                },
            ]
            ({y*cos(x)}, {y*sin(x)}, {sqrt(1-y^2)});

        \end{axis}
    \end{tikzpicture}
\end{center}
\end{eg}

\begin{eg}
    Consider the function $f(x,y) = 2x + 1$. The domain of $f$ is $\R^2$, since $f$ is defined for all points $(x,y) \in \R^2$. The graph of $f$ is a plane in $\R^3$ given by the equation:
    \[
        z = 2x + 1
    \]
    Graphically:
    \begin{center}
        \begin{tikzpicture}
            \begin{axis}[
                axis lines=box,
                axis equal image,
                axis on top,
                z buffer=sort,
                view={35}{20},
                xlabel={$x$},
                ylabel={$y$},
                zlabel={$z$},
                xmin=-5, xmax=5,
                ymin=-5, ymax=5,
                zmin=-10,  zmax=10,
                xtick={-5,0,5},
                ytick={-5,0,5},
                ztick={-10,-5,0,5,10},
                ticks=both,
                enlargelimits=false,
                unit vector ratio=1 2 0.5,
            ]

            \addplot3[
                surf,
                domain=-5:5,
                y domain=-5:5,
                samples=40,
                samples y=20,
                shader=interp,
                colormap={onlyprimary}{
                    color=(primary)
                    color=(secondary)
                },
            ]
            ({x}, {y}, {2*x + 1});

            \end{axis}
        \end{tikzpicture}
    \end{center}
\end{eg}

\begin{eg}
    Consider the function $f(x,y) = ax + by + c$ where $a,b,c \in \R$. If $c = 0$, then we have:
    \begin{align*}
        F & = \{(x,y,z) \in \R^3 : \langle(x,y,z),(a,b,-1)\rangle = 0 \} \\
          & = \{(x,y,z) \in \R^3 : (x,y,z) \perp (a,b,-1)\}
    \end{align*}
    Which is a plane in $\R^3$ passing through the origin and orthogonal to the vector $(a,b,-1)$. If $c \in \R$, then we need to move the plane by $c$ units in the direction of the z-axis.
\end{eg}